\cotefile{103}

\scanpage{15}

Je te signale que j'ai réfléchi aux Gr-champs sur $X$. Si $G$ est un Groupe sur $X$; $N$ un $G$-Module les Gr-champs sur $X$ ``épinglés par $G,N$'' forment a priori une $2$-catégori et même une $2$-catégorie de Picard stricte, grâce à l'opération évident à la Baer. On trouve que le complexe \inserted{(de cochaines)} tronqué à l'échelon $2$ qui leur correspond est le tronqué
\[
\tau_{\leq 2}\bigl(\mathbb R\Gamma(B_{G}\mathbin{\mathrm{mod}}X,N)[1]\bigr)
\]
(\uline{NB} la cohomologie de $\mathbb R\Gamma(B_{G}\mathbin{\mathrm{mod}}X,N)$ commence en degré $1$) Plus géométr\unsure{i}quement, un Gr-champ sur $X$ éping \inserted{par $(G,N)$} est essentiellement ``la même chose'' qu'une $2$-gerbe sur $B_{G}$, liée par $N$, et \deleted{\unsure{épi}}munie d'une trivialisati au dessus de $X\approx B_{\unsure{e|c}}=(B_{G})_{/P}$ (où $P$ est le \deleted{\unsure{``|s|$\ast$}}objet de $B_{G}$ «torseur universel sous $G$»). Ces $2$-gerbes forment en fait une $3$-catégorie de Picard a priori, mais il se
% archive sheet number: 14

\scanpage{16}
trouve que dans celle-ci, les $3$-flèches sont triviales (i.e. si source $=$ but, ce sont des identités) - cela un fait qui exprime $H^{0}(B_{G}/X,N)=0$ (i.e. $H^{0}(B_{G},N)\to H^{0}(X,N)$ injectif\ldots). Donc la $3$-catégorie peut être regardée comme une $2$-catégorie - et ``c'est'' celle des Gr-champs sur $X$ épinglés par $G,N$. Si on veut localiser sur $X$, et décrire le $2$-\uline{champ} \inserted{de Picard} sur $X$ des champs de Picard (sur des objets variables de $X$) épinglés par $G,N$, on trouve qu'il est exprime par le complexe
\[
\tau_{\leq 2}\bigl(\mathbb{R}p_{G*}\operatorname{Coker}\bigl(N\to Rq_{G*}\overbrace{C(q_{G}^{*}N)}^{\text{rés. inj.}}\bigr)\bigr)
\]
où $p_{G}:B_{G}\to X$ et $q_{G}:B_{e}\approx X\simeq(B_{G})_{/\unsure{P|p|\rho}}\to B_{G}$.

Toutes ces descriptions étant compatibles avec des variations de $G,N,X$, cela do\unsure{\deleted{nne}|\deleted{nn}e} en principe une description de la $2$-catégorie des Gr-champs, avec $X,G,N$ variables\ldots
% archive sheet number: 15

\scanpage{17}

\noindent\hfill\unsure{Villecun|Villecui} le 23.6.1974\\

Chers Deligne, Giraud, Verdier,\\

Vous savez peut-être qu'une mathématicienne vietnamienne, Hoang Xuan Sinh, \inserted{a} travaillé depuis mon séjour au Vietnam en 1967 à un sujet de recherche que je lui avais proposé, savoir une théorie de structure des (Gr)-catégories, et la construction d\deleted{\unsure{b|h|n}}'une (Gr)-catégorie (ou plutôt catégorie de Picard) enveloppante pour une catégorie avec produit $\unsure{\otimes|\&}$ associatif et commutatif. J'ai reçu récemment une copie du manuscript en forme (dont la confection s'est faite au prix de mille difficultés - il est douteux qu'il soit possible d'avoir d'autres copies dans un avenir proche, à moins d'en faire d'après cet exemplaire-là), et j'ai envoyé un rapport officiel au sujet de ce travail à M. Ta Quang B$\overset{\handwritten{\unsure{\,)}}}{\mathrm{u}}\handwritten{\breve}{\mathrm{u}}$, ministre de l'enseignement supérieur et technique de la RDV, rapport dont je vous envoie une copie ci-joint. Mme Sinh, avec mon accord, a demandé à présenter ce travail comme thèse de doctorat, et à faire sa soutenance en France. \deleted{\unsure{I|T}}A l'heure où j'écris, je n'ai pas connaissance d'un accord des autorités vietnamiennes pour l'un ni pour l'autre; il y a eu un certain nombre de difficultés soulevées, qu'il est intuile de détailler ici, et qui j'espère vont s'aplanir. J'ai appris par ailleurs que Mme Hoang \unsure{\deleted{H}X|\deleted{X}H}uan Sinh figurerait sur la liste officielle de la délégation vietnamienne au congrès international de mathématiciens à Vancouver du 20 août 1974. Si c'est bien ainsi, il serait très souhaitable qu'elle puisse profiter de cette occasion pour faire sa soutenance à Orsay au retour de Vancouver, fin Septembre ou début Octobre, et il me semble plausible que les autorités vietnamiennes y donneront leur accord. Aussi il serait assez urgent, vu les lenteurs administratives, de constituer un jury. Par la présente, je vous demande si vous seriez disposé en principe à faire partie (avec moi) de ce jury de thèse - ce qui suppose évidemment que vous vous trouvez dans la région parisienne au moment prévu pour la soutenance. A celui qui m'en ferait la demande, je peux envoyer le texte de la thèse complète, ou un résumé assez détaillé rédigé par Mme Hoang Xuan Sinh , ou enfin un résumé moins détaillé de 13 pages manuscrites (rédigées en vue d'une publication aux CR, mais beaucoup trop long à cet effet). Ce serait gentil de me répondre rapidement, et en cas d'impossibilité, de voir autour de vous s'il y a des gens ``dans le bain'' d\deleted{\unsure{e|a}}\handwritten{u} \deleted{\unsure{ce}}\handwritten{genre} \handwritten{\inserted{de choses}} que fait Mme Sinh, et qui seraient intéressés à faire partie du jury.
% archive sheet number: 16

\scanpage{18}

Pour ce qui est des formalités administratives, c'est le frère de Mme Hoan Xuan Han, qui habite à Antony, qui s'en occupera pour elle. A toutes fin utile, je vous passe son adresse:

Hoang Xuan Sanh, 49 rue de Châtenay, Estérel, 92 Antony, Tél BER 6\unsure{8|3} 79.\\

Dans l'atten\unsure{\uline{t}|t}e d'une réponse prochaine, bien cordialement\\
\hfill\handwritten{Schurik}\\

PS N'ayant pas l'adresse de Giraud, je demande à Verdier s'il peut bien lui transmettre la lettre et le rapport. Je,pense que celui-ci\_doit pouvoi\rlap{\raisebox{-1.1ex}[0pt][0pt]{\hspace{.2em}\unsure{(}}}r servir comme rapport de thèse aussi vis à vis de l'administration universitaire en France.
% archive sheet number: 17
